\chapter{Earwax (Cerumen Impaction)}

\section*{Definition}
Earwax (cerumen) is a natural substance produced by ceruminous glands in the outer ear canal that protects the ear by trapping debris and repelling insects. Cerumen impaction occurs when wax accumulates to the point of causing symptoms, affecting approximately 10 percent of children and 5 percent of adults.

\section*{What Happens in Your Body}
Cerumen is composed of shed skin cells, hair, and secretions from ceruminous and sebaceous glands. It normally migrates outward from the tympanic membrane along the ear canal, aided by jaw movement. Impaction occurs when this migration fails due to narrow canals, excessive hair, dry cerumen type, or use of cotton swabs which push wax deeper.

\section*{Causes}
\begin{itemize}
\item Cotton swab use (pushing wax deeper)
\item Anatomical factors (narrow, tortuous ear canals)
\item Excessive ear hair
\item Dry cerumen phenotype (genetic variation)
\item Hearing aid or earplug use
\item Advanced age (drier cerumen)
\item Occupational exposure to dusty environments
\item Excessive swimming or water exposure
\end{itemize}

\section*{When to See a Doctor}
\begin{itemize}
\item Decreased hearing or feeling of fullness in the ear
\item Earache or discomfort
\item Tinnitus (ringing in the ear)
\item Dizziness or vertigo
\item Itching in the ear canal
\item Cough (auricular branch of vagus nerve stimulation)
\item Visible wax blocking the ear canal
\end{itemize}

\section*{Related Symptoms}
\begin{itemize}
\item Hearing loss (conductive)
\item Earache (otalgia)
\item Tinnitus
\item Dizziness or vertigo
\item Otitis externa (swimmer's ear)
\item Ear fullness or pressure
\end{itemize}

\section*{Self-Care / Home Management}
\begin{itemize}
\item Do NOT use cotton swabs or other objects to clean the ear canal
\item Use OTC cerumenolytic drops (carbamide peroxide, mineral oil)
\item Tilt head and allow drops to remain for several minutes
\item Irrigate gently with warm water using a bulb syringe (if no eardrum perforation)
\item Dry the ear canal afterward
\item Never attempt to remove wax if eardrum is perforated
\end{itemize}

\section*{How Your Doctor Will Evaluate You}
\begin{itemize}
\item Otoscopic examination to visualize cerumen and tympanic membrane
\item Tympanometry if hearing loss is present
\item Audiometry if needed
\item Rule out other causes of hearing loss
\end{itemize}

\section*{Treatment Options}
\begin{itemize}
\item Cerumenolytic agents to soften wax
\item Ear irrigation with warm water (after ensuring intact tympanic membrane)
\item Microsuction or manual removal under otoscopic guidance
\item Referral to otolaryngologist for complicated or recurrent impaction
\end{itemize}
\section*{References}
\begin{thebibliography}{99}
\bibitem{ref1} Schwartz SR, Magit AE, Rosenfeld RM, et al. "Clinical Practice Guideline (Update): Earwax (Cerumen Impaction)." Otolaryngol Head Neck Surg. 2017;156(1 Suppl):S1--S29.
\bibitem{ref2} Guest JF, Greener MJ, Robinson AC, et al. "Impacted Cerumen: Composition, Epidemiology, and Management." Clin Otolaryngol Allied Sci. 2004;29(6):570--577.
\bibitem{ref3} Roland PS, Smith TL, Schwartz SR, et al. "Clinical Practice Guideline: Cerumen Impaction." Otolaryngol Head Neck Surg. 2008;139(3 Suppl 2):S1--S21.
\bibitem{ref4} Wright T. "Ear Wax." BMJ. 2015;351:h3601.
\bibitem{ref5} Burton MJ, Doree CJ. "Ear Drops for the Removal of Ear Wax." Cochrane Database Syst Rev. 2018;(7):CD004326.
\end{thebibliography}


\clearpage
